% !TeX spellcheck = nl_NL
\chapter{Algemene \LaTeX{}-tips en -trucs}
Nog enkele algemene \LaTeX{}-tips en -trucs:
\begin{itemize}
    \item Latex kent minstens vier verschillende verschillende liggende streepjes, zie \cref{tab:streepjes}.
        De laatste moet je als rekenkundig minteken gebruiken.
    \tabel{Liggende streepjes in \LaTeX{}}{streepjes}{ccc}{%
        Engelse naam\footnotemark{} & code & resultaat
    }{%
        hyphen & \lcode{-} & - \\
        en dash & \lcode{--} & --  \\
        em dash & \lcode{---} & --- \\
        minus sign & \lcode{\$-\$} & $-$
    }
    \footnotetext{Zie: \url{https://en.wikipedia.org/wiki/Dash}.}
    \item |~| geeft een \enquote{harde} spatie.
    De regel wordt niet afgebroken bij deze harde spatie.
    Voorbeelden:
    \begin{itemize}
       	\item
       	|lengte~$l$|. Het is een raar gezicht als de regel na \dquote{lengte} wordt afgebroken. Het gebruik van de harde spatie voorkomt dit. 
       	De lengte~$l$ wordt opgegeven in \si{\meter}.
       	\item
       	|1, 2 of~3.| 
       	De barman vroeg nogmaals: \dquote{Hoeveel biertjes wil je? 1, 2 of~3?}.
       	\item 
       	|Verhoog de $teller$   met~1 modulo~2.|
       	Verhoog de $teller$ met~1 modulo~2.
    \end{itemize}
    \item |\,| geeft een halve spatie.
    Bij gebruik van en-dash in ranges is bijvoorbeeld een beetje ruimte links en
    rechts nodig. Vergelijk pagina |1--7| met pagina |1\,--\,7|. Vergelijk pagina
    1–-7 met pagina 1\,-–\, 7.
    \item Er is een verschil tussen |\textbf{tekst1}| en |{\bfseries tekst2}|. In |tekst2| mag een paragraafeinde voorkomen, in |tekst1| niet.
    \item package |menukeys|\footnote{\url{https://www.ctan.org/pkg/menukeys}} is handig als je menu-opties en toetsaanslagen wilt benadrukken in tutorials.
    \item Kolomkoppen behoren gecentreerd te zijn.
        De tabelcommando's in de stijl doen dat automatisch.
    \item Als je in math\hyp{}mode een regelmatige rij met getallen wilt opgegeven, dan kan het commando |\ldots| van pas komen: $2, 4, 6, \ldots{}, 100$.
        Als je een regelmatige som van getallen wilt opgegeven, dan geeft het commando |\cdots| een mooier resultaat: $2 + 4 + 6 + \cdots{} + 100$.
    \item Gebruik |\ensuremath| om eigen math\hyp{}commando's in text\hyp{} en math\hyp{}mode te kunnen gebruiken
    \item Bibtex:
    \begin{itemize}
        \item Gebruik |and| tussen namen van verschillende auteurs.
        \item Gebruik accolades rond dubbele achternamen.
        \item Neem \code[mathescape,language={[AlLaTeX]TeX}]{Langid = \{$\ldots{}$\}} op in elke entry zodat er correct wordt afgebroken.
        \item Gebruik |--| bij het aangeven van een aangesloten aantal pagina's:\\ |Pages = {293--326}|.
        \item Geef ISBN nummers op met streepjes ertussen:\\ |ISBN = {978-0-201-89684-8}| om een overfull hbox te voorkomen.
            Gebruik: \url{http://tools.wmflabs.org/isbn/IsbnCheckAndFormat} om ISBN nummers zonder streepjes om te zetten in de officiële notatie (met steepjes).
    \end{itemize}
    \item |\sloppy| zorgt ervoor dat overfull hbox vermeden wordt door spaties ongelimiteerd lang te maken.
        Gebruik daarna |\fussy| om weer netjes uit te vullen.
        In de ebook versie en ook in de papieren versie geeft de url bij het bovenstaande puntje een overfull hbox.
        Kijk wat er met |\sloppy| gebeurt: 
    \sloppy
    \begin{itemize}
        \item
          Gebruik: \url{http://tools.wmflabs.org/isbn/IsbnCheckAndFormat} om ISBN nummers zonder streepjes om te zetten in de officiële notatie (met steepjes).%
    \end{itemize}
    \fussy
    \item Gebruik |\ | tussen een kleine letter en een punt die niet het einde van een zin is.
        B.v.\ bij enz.\ en andere afko's.
        De \LaTeX{}-code van de vorige zin is: |B.v.\ bij enz.\ en andere afko's.|
    \item Gebruik |\@| tussen een hoofdletter en een punt, vraagteken of uitroepteken aan het einde van een zin (om einde zin spatie te krijgen).
        Ik gebruik geen vitamine~C\@.
        Gebruik jij wel vitamine~C\@? Of ook niet.
        De \LaTeX{}-code van de vorige drie zinnen is:
        |Ik gebruik geen vitamine~C\@. Gebruik jij wel vitamine~C\@? Of ook niet.|
    \item Een moving argument is een argument dat ook op een andere plaatst dan waar het gedefinieerd wordt, wordt gebruikt.
        B.v.\ het argument van \hcode{\\sec\\-tion} wordt ook in de inhoudsopgave gebruik.
        Fragile commands mogen niet in moving argumenten gebruikt worden.
        Als je toch een fragile command als moving argument wil gebruiken, dan moet je er |\protect| voorzetten.
        Bijvoorbeeld als je een |\footnote| in een |\section| wilt opnemen:
\begin{lstfragment}[style=lstyle]
\hcode{\section{\texorpdfstring{Titel\protect \footnote{Bedenk zelf een goede titel}}{Titel}}}
\end{lstfragment}
    \section{\texorpdfstring{Titel\protect \footnote{Bedenk zelf een goede titel}}{Titel}}
    \item |\texorpdfstring| is handig als je \LaTeX\hyp{}code in een titel wilt gebruiken.
        De titel komt namelijk ook in de inhoudsopgave van het pdf document maar daar mag geen \LaTeX\hyp{}code in gebruikt worden.
        Bijvoorbeeld:
\begin{lstfragment}[style=lstyle]
section{\texorpdfstring{$E=mc^2$}{E=mc2}}
\end{lstfragment}
    \section{\texorpdfstring{$E=mc^2$}{E=mc2}}
    \item |\centering| start \emph{geen} nieuwe paragraaf \code[mathescape, style=lstyle]{\\begin\{center\} $\ldots$ \\end\{center\}} wel.
    \item Gebruik liever geen absolute lengtes zoals |cm|, |mm| en |pt| maar relatieve zoals |em| (breedte van een M) of |ex| (hoogte van een x).
    \item Het eerste woord in een zin wordt nooit afgebroken, dat kan wel eens problemen geven in een smalle kolom in een tabel.
        Zet er |\hspace{0pt}| voor om het toch af te laten breken.
    \item Gebruik een |\mbox| om ongewenst afbreken te voorkomen.
    \item Bij het maken van een toets met behulp van de class exam kun je een |\solution| op een nieuwe regel beginnen met |\ifprintanswers\\\fi|.
    \item Soms wordt een extra lege regel toegevoegd aan het einde van een paragraaf als de laatste regel van de paragraaf exact past binnen de tekstbreedte. Dit kun je oplossen door |\vskip 0pt| achter de laatste regel van de paragraaf te plaatsen.
    \item \#, \$, \%, \textasciitilde, \_, \textasciicircum, \textbackslash, \{ en \} zijn speciale karakters en moeten als volgt ingevoerd worden: 
    |\#|, |\$|, |\%|, |\textasciitilde|, |\_|, |\textasciicircum|, |\textbackslash|, |\{| en |\}|.
    \item In lopende tekst kun je het commando |\textcite| gebruiken om de naam van de auteur met de verwijzing te genereren. De code |\textcite{latex} beweert \ldots{}| levert de volgende tekst op: \textcite{latex} beweert \ldots{}.
    \item Een leesteken tussen \textbf{vet}, normaal\emph{, cursief} moet in de lichtste font worden opgemaakt.
    \item Gebruik smallcaps voor afkortingen zoals |\textsc{html}|: \textsc{html}.
    Zo ziet \textsc{ip}-adres er mooier uit dan IP-adres.
    \item Het woord finish wordt door \LaTeX{} opgemaakt met een zogenaamde ligatuur\footnote{%
    	Zie: \url{https://nl.wikipedia.org/wiki/Ligatuur_(typografie)}.
    }. In woorden waarbij de f en de i niet bij dezelfde lettergreep horen, zoals bijvoorbeeld selfish is het beter om  om geen ligatuur te gebruiken.
    Gebruik daarom \code{self\\makebox\{\}ish} self\makebox{}ish.
    Vergelijk selfish met self\makebox{}ish\footnote{%
    	Je moet wel inzoemen om het verschil te kunnen zien. In selfish zit de f vast aan de i en bij het selecteren van tekst is de fi één letter. In self\makebox{}ish zit een beetje ruimte tussen de f en de i en bij het selecteren van tekst zijn het twee aparte letters.  
	}.
    Hetzelfde geldt voor woorden zoals off\makebox{}line enzovoorts.
    \item Gebruik de package |upgreek| voor rechtopstaande griekse letters.
\end{itemize}

Nog enkele tips voor het gebruik van \LaTeX{} in combinatie met git:
\begin{itemize}
    \item 
	    Als de \LaTeX{}-code in een VCS wordt opgenomen worden alle wijzigingen en contributies netjes bijgehouden.
        Binnen de opleiding Elektrotechniek van de Hogeschool Rotterdam maken we gebruik van BitBucket (zie: \url{https://bitbucket.org/HR_ELEKTRO/}).
        We maken gebruik van BitBucket en niet van GitHub omdat de eerste gratis private repositories biedt.
    \item 
	    Het is handig om na elke zin een enter in te typen.
	    Git vergelijkt bestanden namelijk regel voor regel en wijzigingen worden dus veel duidelijker weergegeven als de regels niet te lang zin.
\end{itemize}
